% File: Volume-I-Foundations/chapters/ch02-infinity-categories.tex

\chapter{\texorpdfstring{$\infty$}{Infinity}--Categories and Homotopy Theory}
\label{chap:inf-cat}

This chapter reviews the homotopical and $\infty$--categorical structures
needed to formulate derived geometry and, ultimately, the lamphrodynamic moduli
stack $\Plenum$ developed in later chapters. We assume that the reader is
familiar with basic category theory and classical homotopy theory. Our goal is
not to give a comprehensive treatment, but to present a minimal collection of
definitions, constructions, and results that will be invoked in the sequel.

Throughout, we work over a field of characteristic zero (typically $\R$ or
$\C$) unless otherwise specified. References include \cite{Lurie2009HTT},
\cite{Lurie2017HA}, and \cite{ToenVezzosi2008}.

%---------------------------------------------------------------
\section{Model Categories and Homotopical Structures}
%---------------------------------------------------------------

The motivation for $\infty$--categories begins with the desire to regard
homotopy equivalent objects as ``the same'' while retaining control over
higher homotopies. Classical category theory does not remember higher
homotopies; instead, it collapses homotopies to identities. In homotopy
theory, however, we need to keep track of:
\begin{itemize}
  \item homotopies between maps,
  \item homotopies between homotopies,
  \item and so on, ad infinitum.
\end{itemize}

A traditional approach is provided by \emph{model categories}, which endow a
category with distinguished classes of weak equivalences, fibrations, and
cofibrations satisfying Quillen's axioms \cite{Quillen1967}. The basic example
is the model structure on topological spaces, where weak equivalences are
weak homotopy equivalences, fibrations are Serre fibrations, and cofibrations
are retracts of relative CW--complex inclusions.

While model categories provide a powerful toolkit, they are ultimately an
\emph{encoding} of $\infty$--categorical structures, and the language of
$\infty$--categories gives a more conceptual, coordinate--free approach.

%---------------------------------------------------------------
\section{\texorpdfstring{$\infty$}{Infinity}--Categories}
%---------------------------------------------------------------

An $\infty$--category is a category enriched in spaces rather than sets, or,
equivalently, a homotopical object that remembers all higher morphisms. A
concrete model is given by \emph{quasi--categories} (simplicial sets
satisfying inner horn conditions) introduced by Boardman and Vogt and
developed extensively by Joyal and Lurie.

\begin{definition}
A \emph{quasi--category} is a simplicial set $C$ such that every inner
horn $\Lambda^n_i \to C$ admits a filler $\Delta^n\to C$ for $0 < i < n$.
\end{definition}

This condition generalizes the Kan condition for simplicial sets
corresponding to $\infty$--groupoids (spaces). Every topological space has a
singular complex which is a Kan complex; every Kan complex is a model for an
$\infty$--groupoid. The quasi--category condition weakens Kan filling so that
objects need not have all inverses, thus modeling $\infty$--categories rather
than $\infty$--groupoids.

\Intuition
Ordinary categories forget higher homotopies between morphisms. Quasi--categories
remember not only morphisms but also homotopies and homotopies between homotopies,
all the way up. This makes them ideal for encoding moduli spaces, where higher
automorphisms naturally appear.

%---------------------------------------------------------------
\section{$\infty$--Topoi and Sheaves of Spaces}
%---------------------------------------------------------------

Derived geometry requires a setting where ``sheaves of homotopical data'' can
be treated systematically. This motivates the notion of an $\infty$--topos:
a homotopical enhancement of Grothendieck topoi.

\begin{definition}
An \emph{$\infty$--topos} is an $\infty$--category that is equivalent to an
$\infty$--category of sheaves of spaces on a Grothendieck site. Equivalently,
it is a presentable $\infty$--category satisfying descent for hypercovers
\cite{Lurie2009HTT}.
\end{definition}

Classical sheaf theory associates sets or vector spaces to open subsets of a
manifold. An $\infty$--topos associates \emph{spaces} (i.e.\ $\infty$--groupoids)
to open subsets, thereby encoding higher homotopical information. Derived
stacks are then defined as (higher) sheaves of $\infty$--groupoids satisfying
appropriate representability and local geometricity conditions.

\begin{remark}
The $\infty$--topos of sheaves of spaces on the site $(\mathrm{Sch},\mathrm{Zariski})$
is denoted ${\rm Shv}_{\infty}(\mathrm{Sch})$. Derived algebraic geometry works
mostly inside suitable subcategories of this $\infty$--topos.
\end{remark}

%---------------------------------------------------------------
\section{Derived Stacks}
%---------------------------------------------------------------

Derived stacks were developed by Toën–Vezzosi, Lurie, and others to handle
moduli problems with obstructions, singularities, and nontrivial automorphism
groups. Roughly speaking, derived stacks are functors from derived rings to
$\infty$--groupoids satisfying descent and representability conditions.

\begin{definition}
A \emph{derived stack} is a functor
\[
  F : \mathrm{dAlg}^{op} \longrightarrow \infty\text{--}\mathrm{Gpd}
\]
satisfying $\infty$--geometricity and descent conditions.
\end{definition}

Derived stacks generalize classical schemes and algebraic stacks, adding:
\begin{itemize}
  \item derived intersections,
  \item cotangent complexes,
  \item obstruction theories,
  \item and higher automorphisms.
\end{itemize}

\Intuition
In physics, moduli spaces of solutions to gauge equations (e.g.\ Yang--Mills,
GR) often have singularities and higher gauge symmetries. Derived stacks treat
these automatically, removing the need for ad hoc patching arguments.

In \rsvpabbr{}, the moduli stack $\Plenum$ of lamphrodynamic configurations
will be constructed as a derived Artin stack equipped with a $(-1)$--shifted
symplectic structure. This stack encodes field configurations, gauge orbits,
and entropic singularities in a unified homotopical object.

%---------------------------------------------------------------
\section{Tangent Complexes and Obstruction Theory}
%---------------------------------------------------------------

Classical manifolds possess tangent spaces at each point, controlling
infinitesimal deformations. Derived stacks generalize this to a \emph{tangent
complex}, which is typically a chain complex, possibly with cohomology in
multiple degrees.

\begin{definition}
Let $x\in F$ be a point of a derived stack $F$. The \emph{tangent complex}
$T_xF$ is the derived mapping space
\[
  T_xF := \Map_{\mathrm{dStk}}(\Spec k[\epsilon]/(\epsilon^2), F)_x,
\]
considered as a complex in degrees $(-\infty,\dots,\infty)$.
\end{definition}

The cohomology groups of $T_xF$ encode:
\begin{itemize}
  \item $H^0(T_xF)$ classical tangent directions,
  \item $H^{-1}(T_xF)$ infinitesimal automorphisms,
  \item $H^1(T_xF)$ obstructions to extending deformations,
  \item higher degrees controlling higher obstructions and symmetries.
\end{itemize}

In particular, if $F$ is smooth (in derived sense), then $H^1(T_xF)$ may vanish,
ensuring unobstructed deformations; if $H^{-1}(T_xF)$ is nonzero, $x$ has
nontrivial automorphisms (gauge symmetries).

\Intuition
The tangent complex of $\Plenum$ will encode:
\begin{itemize}
  \item infinitesimal variations of $(\PhiField,\vField,\SField)$,
  \item gauge symmetries mixing $\PhiField$ and $\vField$,
  \item obstructions related to entropic singularities.
\end{itemize}

%---------------------------------------------------------------
\section{Cotangent Complex and Shifted Symplectic Structures}
%---------------------------------------------------------------

The cotangent complex $L_F$ of a derived stack $F$ is a fundamental invariant
generalizing the classical cotangent bundle. It plays a central role in the
definition of shifted symplectic structures.

\begin{definition}
The \emph{cotangent complex} $L_F$ is a perfect complex on $F$ (when $F$ is
derived Artin) whose dual $T_F := L_F^\vee$ generalizes the tangent bundle.
\end{definition}

Shifted symplectic structures are defined as certain closed bilinear forms on
$L_F$ (or $T_F$) of cohomological degree $k$. The case $k=-1$ is the setting of
the BV formalism and will be central to \rsvpabbr{}.

We discuss shifted symplectic structures in detail in
\cref{chap:shifted-symplectic}.

%---------------------------------------------------------------
\section{\texorpdfstring{$\infty$}{Infinity}--Groupoids and Fields}
%---------------------------------------------------------------

In ordinary geometry, a point in a moduli space of fields corresponds to a
single classical solution. In derived geometry, such points carry higher
homotopical structure: their automorphism groups form $\infty$--groupoids,
which encode gauge equivalences and higher symmetries.

Consequently, the moduli stack $\Plenum$ should be thought of as an object in
an $\infty$--category of derived stacks, with points representing field
configurations, morphisms representing gauge transformations, and higher
morphisms encoding redundancy among transformations themselves.

This viewpoint is essential to the BV formalism: ghosts, antifields, and the
master equation all arise naturally from the derived structure of $\Plenum$.

%---------------------------------------------------------------
\section{Summary and Outlook}
%---------------------------------------------------------------

This chapter introduced the basic machinery underlying derived geometry:
model categories, $\infty$--categories, $\infty$--topoi, and derived stacks,
with an emphasis on tangent and cotangent complexes and higher automorphisms.
These tools will be used in \cref{chap:shifted-symplectic,chap:aksz-bv} to
define and analyze shifted symplectic structures and the AKSZ formalism.

In \cref{chap:plenum-stack}, we will construct the lamphrodynamic moduli
stack $\Plenum$ explicitly as a derived Artin stack, then equip it with a
$(-1)$--shifted symplectic form in
\cref{chap:aksz-rsvp}. The present chapter thus serves as the homotopical
foundation for all subsequent developments.


